%%%%%%%%%%%%%%%%%%%%%%%%%%%%%%%%%%%%%%%%%%%%%%%%%%%%%%%%%%%%%%%%%%%%%%%%%%%%
% CAPÍTULO 10: ARQUITECTURA Y SOLUCIÓN PROPUESTA
%%%%%%%%%%%%%%%%%%%%%%%%%%%%%%%%%%%%%%%%%%%%%%%%%%%%%%%%%%%%%%%%%%%%%%%%%%%%

\chapter{Arquitectura y solución propuesta}
\label{ch:arquitectura}

\section{Visión general (conceptual)}
\label{sec:vision-conceptual}

La solución se estructura en seis bloques funcionales de alto nivel, que separan responsabilidades y permiten sustituir componentes sin rehacer el sistema:
\begin{enumerate}[label=\alph*)]
    \item \textbf{Bloque de Presentación (\textit{Frontend}):} Punto de acceso para conductores y pasajeros. Se implementa con Angular + Ionic, lo que permite una única base de código que se compila para web y móvil (RNF-02). Incluye las pantallas de publicación de viajes, búsqueda, solicitudes y visualización de rutas.
    \item \textbf{Bloque de Administración (Panel Admin):} Interfaz separada para el administrador. Desde aquí se lanzan las simulaciones (escenarios S0--S5), se consultan las métricas y se exportan los resultados en CSV (RF-15 a RF-18, RF-20). Se implementa exclusivamente en web, ya que las tareas de gestión resultan más cómodas en un navegador.
    \item \textbf{Bloque de Comunicación (API Gateway):} Expone los servicios del Core y del Componente de Experimentación al \textit{frontend}. Gestiona autenticación, validación y enrutamiento de peticiones (Node.js).
    \item \textbf{Bloque del Core del Sistema:} Contiene la lógica de negocio: gestión de usuarios, viajes y solicitudes; el algoritmo de \textit{matching}; y el Gemelo Digital. Este mantiene una representación sincronizada de usuarios, vehículos, viajes, solicitudes y rutas, y se actualiza mediante eventos.
    \item \textbf{Bloque de Experimentación:} Se divide en A (Generador de Escenarios), que crea datos sintéticos y define los parámetros de cada escenario (S0--S5), y B (Motor de Simulación). El Motor B ejecuta el bucle de simulación sobre una copia del estado del Gemelo Digital, sin afectar a los datos reales; aplica el \textit{matching}, asigna plazas y calcula métricas como vehículos evitados, ocupación media y solicitudes cubiertas.
    \item \textbf{Bloque de Persistencia:} Almacena y recupera datos de usuarios, viajes, solicitudes y resultados de simulación. Se implementa con PostgreSQL + PostGIS, que permite realizar consultas geoespaciales para el \textit{matching}.
\end{enumerate}

La Figura~\ref{fig:bloques-funcionales} muestra el diagrama de bloques funcionales de alto nivel. En ella se aprecia la separación entre los bloques funcionales de la plataforma (a, b, c, d, f) y el bloque de experimentación (e), que opera sobre una copia del estado del Gemelo Digital.

\begin{figure}[H]
  \centering
  \begin{tikzpicture}[x=1cm,y=1cm]
    % --- fila superior: los dos clientes (van en paralelo) ---
    \node[bloque, minimum width=7.2cm, minimum height=1.85cm] (a) at (3.95,12.35)
      {\textbf{a) Bloque de Presentación}\\[1pt]
       \textit{Frontend} de usuarios\\[1pt]
       {\scriptsize Angular + Ionic $\rightarrow$ web y móvil}};
    \node[bloque, minimum width=7.0cm, minimum height=1.85cm] (b) at (11.85,12.35)
      {\textbf{b) Bloque de Administración}\\[1pt]
       Panel de administración (web)\\[1pt]
       {\scriptsize lanzamiento de escenarios y consulta de métricas}};

    % --- pasarela de comunicación ---
    \node[bloque, minimum width=14.9cm, minimum height=1.45cm] (c) at (7.9,9.85)
      {\textbf{c) Bloque de Comunicación} --- API Gateway (Node.js)\\[1pt]
       {\scriptsize autenticación, validación y enrutamiento de peticiones}};

    % --- core del sistema ---
    \draw[draw=Color2, fill=Color2!6, rounded corners=2pt] (0.45,4.85) rectangle (9.45,8.05);
    \node[rotulo] at (0.60,7.90) {d) Bloque del Core del Sistema};
    \foreach \x/\t in {1.70/{Gestión de\\Usuarios}, 3.85/{Gestión de\\Viajes},
                       6.00/{\textit{Matching}}, 8.15/{Gemelo\\Digital}} {
      \node[subbloque, minimum width=1.90cm] at (\x,6.05) {\t};
    }

    % --- experimentación ---
    \draw[draw=Color2, fill=Color2!6, rounded corners=2pt] (10.35,4.85) rectangle (15.35,8.05);
    \node[rotulo] at (10.50,7.90) {e) Bloque de Experimentación};
    \node[subbloque, minimum width=2.30cm] at (11.62,6.30) {A) Generador\\de escenarios};
    \node[subbloque, minimum width=2.30cm] at (14.08,6.30) {B) Motor de\\simulación};
    \node[font=\scriptsize, align=center, text width=4.6cm] at (12.85,5.35)
      {opera sobre una \textbf{copia} del estado del Gemelo Digital};

    % --- persistencia ---
    \node[bloque, minimum width=14.9cm, minimum height=1.35cm] (f) at (7.9,2.95)
      {\textbf{f) Bloque de Persistencia} --- PostgreSQL + PostGIS\\[1pt]
       {\scriptsize usuarios, vehículos, viajes, solicitudes y resultados de simulación}};

    % --- conectores ---
    \draw[flecha] (a.south) -- (a.south |- c.north);
    \draw[flecha] (b.south) -- (b.south |- c.north);
    \draw[flecha] (4.60,9.12) -- (4.60,8.05);
    \draw[flecha] (12.60,9.12) -- (12.60,8.05);
    \draw[flecha] (4.60,4.85) -- (4.60,3.63);
    \draw[flecha] (12.60,4.85) -- (12.60,3.63);
  \end{tikzpicture}
  \caption[Diagrama de bloques funcionales]{Diagrama de bloques funcionales de alto nivel. Los bloques de presentación y administración son clientes independientes que entran por la misma API; el bloque de experimentación no accede al estado real, sino a una copia del Gemelo Digital.}
  \label{fig:bloques-funcionales}
\end{figure}

\section{Vista técnica (arquitectura por capas)}

Mientras que la sección anterior presenta la visión conceptual del sistema (bloques funcionales: qué hace cada parte), esta sección presenta la visión técnica (capas: cómo se implementa con tecnologías concretas). Ambas vistas son complementarias y responden al enfoque de múltiples vistas arquitectónicas \cite{Kruchten1995}.

La Tabla~\ref{tab:capas} muestra la arquitectura técnica del sistema.

\begin{table}[H]
  \centering
  \begin{tabularx}{\textwidth}{Y{0.5769} Y{1.5000} Y{0.9231}}
    \hline
    \textbf{Capa} & \textbf{Responsabilidad} & \textbf{Tecnología prevista} \\
    \hline
    Presentación & \textit{Frontend} usuarios (web + móvil) + Panel Admin (web) & Angular + Ionic \\
    API / \textit{Backend} & Autenticación, usuarios, viajes, \textit{matching}, reglas & Node.js \\
    Persistencia & Usuarios, vehículos, viajes, solicitudes, estados, datos de simulación & PostgreSQL + PostGIS \\
    Gemelo Digital & Representación de entidades y estado del sistema & Servicio/módulo DT \\
    Simulación & Generación de escenarios y ejecución de experimentos & Módulo de simulación \\
    DevOps & Control de versiones, pruebas y despliegue & GitHub + GitHub Actions \\
    \hline
  \end{tabularx}
  \caption{Arquitectura por capas.}
  \label{tab:capas}
\end{table}

\textbf{Descripción de las capas:}
\begin{itemize}
    \item \textbf{Capa de Presentación:} Interfaz de usuario para conductores y pasajeros (web + móvil) y panel de administración (web). Se implementa con Angular + Ionic.
    \item \textbf{Capa de API/\textit{Backend}:} Expone los servicios del sistema mediante APIs REST. Se implementa con Node.js.
    \item \textbf{Capa de Persistencia:} Almacena y recupera datos. PostgreSQL + PostGIS permite consultas geoespaciales para el \textit{matching}.
    \item \textbf{Capa de Gemelo Digital:} Mantiene una representación sincronizada del estado del sistema. Se actualiza por eventos.
    \item \textbf{Capa de Simulación:} Genera escenarios (S0--S5) y ejecuta simulaciones sobre una copia del estado del Gemelo Digital.
    \item \textbf{Capa de DevOps:} Capa transversal de soporte al desarrollo y despliegue.
\end{itemize}

\textit{Nota sobre DevOps:} DevOps se incluye en esta vista técnica como capa transversal de soporte al desarrollo y despliegue. No aparece en la vista conceptual (Sección~\ref{sec:vision-conceptual}) porque no constituye un bloque funcional del sistema, sino una infraestructura de soporte que afecta a todo el ciclo de vida.

\subsection{Correspondencia entre bloques funcionales y capas técnicas}

La Tabla~\ref{tab:correspondencia} muestra la correspondencia entre los bloques funcionales y las capas técnicas, reforzando la coherencia entre ambas vistas arquitectónicas.

\begin{table}[H]
  \centering
  \begin{tabularx}{\textwidth}{Y{0.8333} Y{1.1667}}
    \hline
    \textbf{Bloque funcional} & \textbf{Capa técnica} \\
    \hline
    Bloque de Presentación (\textit{Frontend}) & Capa de Presentación \\
    Bloque de Administración (Panel Admin) & Capa de Presentación \\
    Bloque de Comunicación (API Gateway) & Capa de API/\textit{Backend} \\
    Bloque del Core del Sistema & Capa de API/\textit{Backend} + Capa de Gemelo Digital \\
    Bloque de Experimentación & Capa de Simulación \\
    Bloque de Persistencia & Capa de Persistencia \\
    (transversal) & Capa de DevOps \\
    \hline
  \end{tabularx}
  \caption{Correspondencia entre bloques funcionales y capas técnicas.}
  \label{tab:correspondencia}
\end{table}

\section{Modelo del Gemelo Digital}

El Gemelo Digital será una representación computacional del sistema de movilidad compartida gestionado por la plataforma. Se prioriza el nivel de abstracción necesario para responder a la pregunta de investigación, sin modelar elementos físicos complejos del vehículo.

\subsection{Entidades digitales}

Las entidades digitales del Gemelo Digital se detallan en la Tabla~\ref{tab:entidades}.

\begin{table}[H]
  \centering
  \begin{tabularx}{\textwidth}{Y{0.6000} Y{1.2000} Y{1.2000}}
    \hline
    \textbf{Entidad} & \textbf{Atributos principales} & \textbf{Eventos / cambios} \\
    \hline
    Usuario & id, rol, preferencias, origen/destino del viaje & registro, solicitud, aceptación, cancelación \\
    Vehículo & id, capacidad, plazas disponibles, estado, localización & alta, disponibilidad, ocupación, movimiento \\
    Viaje & id, origen, destino, hora, ruta, estado & creación, confirmación, inicio, finalización \\
    Solicitud & pasajero, viaje, estado & creación, aceptación, rechazo, cancelación \\
    Ruta & origen, destino, distancia, duración, secuencia & recálculo o actualización \\
    Escenario & usuarios, vehículos, demanda, reglas & creación, ejecución, comparación \\
    \hline
  \end{tabularx}
  \caption{Entidades digitales del Gemelo Digital.}
  \label{tab:entidades}
\end{table}

El tipo \texttt{Coordenada} se modela como un \textit{value object} reutilizable para representar ubicaciones geográficas. El tipo \texttt{Evento} desempeña un papel clave en la sincronización con el Gemelo Digital: cada cambio de estado de una entidad genera un evento que permite actualizar su representación digital.

La Figura~\ref{fig:clases} muestra el diagrama de clases del dominio.

\begin{figure}[H]
  \centering
  \includegraphics[width=0.8\textwidth]{clases-dominio}
  \caption{Diagrama de clases del dominio del Gemelo Digital.}
  \label{fig:clases}
\end{figure}

\subsection{Principios de sincronización}

\begin{itemize}
    \item Cada entidad relevante del sistema real tiene una representación digital identificable.
    \item Los cambios de estado de la plataforma se reflejan en el modelo digital.
    \item El modelo digital puede consultar el estado actual y utilizarlo como entrada para simulación.
    \item La comunicación entre plataforma y Gemelo Digital se diseña mediante interfaces bien definidas (eventos + API).
\end{itemize}

\subsection{Nivel de fidelidad}

El DT es de baja fidelidad:
\begin{itemize}
    \item Representa estado y eventos, no dinámica física.
    \item No incorpora sensórica ni V2X.
    \item Se alimenta de los eventos generados por la propia aplicación.
\end{itemize}

\section{Lógica de \textit{matching}}

El \textit{matching} constituye una funcionalidad central. Su primera versión se plantea como un algoritmo interpretable basado en compatibilidad de atributos.

\subsection{Criterios de compatibilidad}

Los criterios de compatibilidad se detallan en la Tabla~\ref{tab:criterios}.

\begin{table}[H]
  \centering
  \begin{tabularx}{\textwidth}{Y{0.6923} Y{1.6154} Y{0.6923}}
    \hline
    \textbf{Criterio} & \textbf{Ejemplo} & \textbf{Prioridad} \\
    \hline
    Origen compatible & Distancia máxima entre puntos de origen & Alta \\
    Destino compatible & Destino igual o dentro de una zona definida & Alta \\
    Ventana temporal & Hora de salida $\pm$ tolerancia & Alta \\
    Capacidad & Plazas disponibles del vehículo & Alta \\
    Desvío & Incremento de recorrido por recoger pasajeros & Media \\
    Preferencias & Preferencias declaradas por usuarios & Media \\
    Ocupación & Maximizar número de plazas utilizadas & Media \\
    \hline
  \end{tabularx}
  \caption{Criterios de compatibilidad del \textit{matching}.}
  \label{tab:criterios}
\end{table}

\subsection{Proceso progresivo}

\begin{itemize}
    \item Filtrado geoespacial (PostGIS): eliminar candidatos claramente incompatibles por distancia.
    \item Comprobación temporal: verificar ventana horaria y tolerancia.
    \item Evaluación de ruta: calcular desvío y tiempo adicional.
    \item Puntuación de compatibilidad: ordenar los viajes resultantes.
\end{itemize}

En una fase posterior, si los resultados lo justifican, podrá incorporarse una función de puntuación ponderada o un algoritmo de optimización. Esta ampliación no es necesaria para alcanzar el objetivo mínimo del TFM.

\section{Módulo de simulación}

La simulación es el mecanismo que convierte al Gemelo Digital en una herramienta de evaluación y no solo de representación. Se propone comparar escenarios con diferentes niveles de demanda y participación en el coche compartido.

\subsection{Escenarios previstos}

Los escenarios previstos se detallan en la Tabla~\ref{tab:escenarios}.

\begin{table}[H]
  \centering
  \begin{tabularx}{\textwidth}{Y{0.6429} Y{1.5000} Y{0.8571}}
    \hline
    \textbf{Escenario} & \textbf{Descripción} & \textbf{Indicadores} \\
    \hline
    S0 --- Base & Cada usuario realiza el desplazamiento de forma individual & N.º vehículos, ocupación, viajes \\
    S1 --- Carpooling básico & \textit{Matching} con reglas de compatibilidad & vehículos evitados, ocupación media \\
    S2 --- Mayor demanda & Aumenta el número de usuarios & escalabilidad, tiempos de \textit{matching} \\
    S3 --- Horarios variables & Cambian las ventanas temporales & compatibilidad, viajes cubiertos \\
    S4 --- Capacidad variable & Vehículos con diferente número de plazas & ocupación y eficiencia \\
    S5 --- Incidencias & Cancelaciones o cambios de disponibilidad & robustez y recuperación \\
    \hline
  \end{tabularx}
  \caption{Escenarios previstos.}
  \label{tab:escenarios}
\end{table}

\subsection{Indicadores iniciales}

\begin{itemize}
    \item Número total de viajes solicitados.
    \item Número de viajes compatibles encontrados.
    \item Número de vehículos utilizados.
    \item Número de vehículos potencialmente evitados respecto al escenario base.
    \item Ocupación media de los vehículos.
    \item Porcentaje de solicitudes atendidas.
    \item Tiempo de respuesta del \textit{matching}.
    \item Distancia adicional por recogida, si se modela.
    \item Número de cancelaciones/incidencias simuladas.
\end{itemize}

\section{Trazabilidad arquitectura \texorpdfstring{$\leftrightarrow$}{<->} requisitos \texorpdfstring{$\leftrightarrow$}{<->} objetivos}

La Tabla~\ref{tab:trazabilidad-arq} muestra la trazabilidad entre componentes arquitectónicos, requisitos y objetivos.

\begin{table}[H]
  \centering
  \begin{tabularx}{\textwidth}{Y{1.0800} Y{0.9600} Y{0.9600}}
    \hline
    \textbf{Componente arquitectónico} & \textbf{Requisitos asociados} & \textbf{Objetivos} \\
    \hline
    \textit{Frontend} (Angular + Ionic) & RF-01 a RF-12, RF-19 & OE9 \\
    \textit{Backend} (Node.js) & RF-01 a RF-14, RF-18 & OE7, OE8, OE9 \\
    Persistencia (PostgreSQL/PostGIS) & RF-03 a RF-14 & OE4, OE9 \\
    Gemelo Digital & RF-13, RF-14, RNF-13 & OE4, OE5, OE6 \\
    \textit{Matching} & RF-06, RF-07 & OE7 \\
    Simulación & RF-15 a RF-17 & OE10 \\
    Métricas y panel & RF-18, RF-20 & OE11, OE12 \\
    DevOps (GitHub + Actions) & RNF-05 & OE9 \\
    \hline
  \end{tabularx}
  \caption{Trazabilidad arquitectura \texorpdfstring{$\leftrightarrow$}{<->} requisitos \texorpdfstring{$\leftrightarrow$}{<->} objetivos.}
  \label{tab:trazabilidad-arq}
\end{table}
